\documentclass[10pt,norsk,a4paper,usenames,dvipsnames]{article}
\usepackage[utf8]{inputenc}
\usepackage[T1]{fontenc}
\usepackage[norsk]{babel}
% - PDF-relatert
\usepackage{hyperref,pdfpages,hypcap,download}
\hypersetup{colorlinks=true,allcolors=.}
\newcommand\fhref[2]{%
	\href{#1}{#2}\footnote{\url{#1}}%
}
% Andre pakker
\usepackage[cm]{fullpage}
\usepackage{parskip,multicol,textcomp,amssymb,graphicx,color,enumitem,cleveref}
% - Korreksjon av fotnoter i seksjoner/overskrifter
\usepackage[stable]{footmisc}
% - Skrifttype
\usepackage[bitstream-charter]{mathdesign}
% - Kommentarer
\usepackage{comment}
\usepackage{epstopdf}
\usepackage[dvipsnames]{xcolor}
\usepackage{xfrac}
% Redefines the first level of enumerated lists.
\renewcommand{\theenumi}{\alph{enumi}}
\renewcommand{\labelenumi}{\theenumi)}

\epstopdfsetup{outdir=./}

\download[vedtekter.pdf]{protokoll/cyb-gf-protokoll-2021-vår.pdf}


\title{\huge Protokoll - Ekstraordinær generalforsamling}
\author{\LARGE Cybernetisk Selskab}
\date{27. april 2022}

\begin{document}


\maketitle



\begin{center}


\includegraphics[width=0.6\textwidth,height=0.6\textheight,keepaspectratio=true]{cyblogoa3.pdf}

\end{center}


\newpage


\tableofcontents

\section{Valg av møteleder}

\textbf{Markus Frislid ble valgt som møteleder ved akklamasjon.}

\section{Valg av referent}

\textbf{Markus Christopher Barkenæs ble valgt som referent ved akklamasjon.}

\section{Valg av protokollunderskrivere}

\textbf{Bendik Berg og Fredrik Aas Andreasen ble valgt som protokollunderskrivere ved akklamasjon.}

\section{Valg av tellekorps}

\textbf{Kamilla Lund og Ebba Olson ble valgt som tellekorps ved akklamasjon.}

\section{Godkjenning av innkalling}

\textbf{Innkalling ble sendt ut 19. april. Innkalling ble godkjent ved akklamasjon.}

\section{Godkjenning av dagsorden}

\textbf{Dagsorden ble godkjent ved akklamasjon.}

Møteregler ble presentert av møteleder og godkjent ved akklamasjon.
Formalia ble ferdig 17:50.
\section{Valg}

    \subsection{Hovedstyret}
        Man velges inn i hovedstyret for to (2) semestere av gangen.

    
    \subsubsection{Rekrutteringsansvarlig} 
    \textbf{Adrian Skansen ble valgt ved akklamasjon.}
    

\section{Vedtektsforslag}

Salen er innstilt for å ta debatten, med den tid det trengs for å ta den.

Styreleder presenterer hovedstyrets innstilling til vedtektsdebatten, og innstiller på ikke noen vedtektsendringer. Møteleder presenterer foreslåtte vedtektsendringer. Det ble innledet spørsmålsrunde. Debatten og votering ble tatt forslag for forslag, etter debatt hvorvidt det burde gjøres slik eller om voteringen burde skje samlet på slutten. Vi har 78 stemmeberettigete i salen. 40 trengs for alminnelig flertall. 52 for kvalifisert flertall. 

    \subsection{§5f - etterfylling}
        Endre:
        \\Den etterfylte sitter frem til neste generalforsamling.
        \\Til:
        \\Den etterfylte sitter frem til neste \textcolor{ForestGreen}{ordinære} generalforsamling.
        
        \textbf{Forslaget er vedtatt etter kort debatt.}
        
    \subsection{§6e - etterfylling}
        Endre:
        \\Hovedstyret kan oppnevne kjellerstyremedlemmer ved behov for etterfylling av alle verv foruten kjellermogul.
        \\Til:
        \\Hovedstyret kan oppnevne kjellerstyremedlemmer ved behov for etterfylling av alle verv foruten kjellermogul. \textcolor{ForestGreen}{Den etterfylte sitter frem til generalforsamling.}
        
        \textbf{Forslaget falt. Kort debatt med mange avholdene.}

    \subsection{§7g og 7e}
    
        \small\textit{(Forslaget denne endringen er basert på er blitt trukket)}\\\normalsize
        Endre:
        \\Bytte plass på 7e og 7g
        
        \textbf{Forslaget endringen er basert på er trukket, enighet om å ikke behandle.}
    
    \subsection{§1 - Formål}
        Legge til nytt punkt (c) i paragraf 1:
        \\a) Cybernetisk Selskab er en instituttforening ved Institutt for informatikk ved Universitetet i Oslo.
        \\b) Cybernetisk Selskabs formål er å arrangere og fremme aktiviteter og arrangementer for studenter ved instituttet. Gjennom foreningens arbeid ønsker man å fremme faglig innhold, skape kameratslig samvær og bidra til et trygt miljø i studietiden.
        \\\textcolor{ForestGreen}{c) Cybernetisk Selskab skal etterstrebe å driftes av aktive studenter. Dette gjelder, men er ikke begrenset til verv, kafédrift, intern kommunikasjon m.m.}
        
        \textbf{Forslaget ble enstemming vedtatt med unntak av 2 avholdene. Bred enighet i salen.}
    
    \subsection{§2 - Medlemskap}
        Endre teksten i paragraf 2 (a), samt fjerne punkt (b) og (c) og flytte punkt (d) til punkt (b)
        \\\\\textbf{Slik den er nå:}
        \begin{enumerate} 
        	\item{Cybernetisk Selskab har tre kategorier medlemmer:}
        	\begin{itemize}
        		\item{Semesterbetalende medlemmer.}
        		\item{Livsvarige medlemmer.}
        		\item{Æresmedlemmer.}
        	\end{itemize}
        	\item{Kontingent for livsvarlig medlemskap er ti ganger ordinær semesterkontingent.}
        	\item{Æresmedlemmer utnevnes på generalforsamling med \sfrac{2}{3} flertall.}
        	\item Medlemmer forplikter seg til å handle i tråd med foreningens formål jf. §1 og etiske retningslinjer for forenings- og styremedlemmer ved instituttet.\\
        \end{enumerate}
        
        \textbf{Foreslått til:}
        \begin{enumerate} 
        	\item{Medlemskap i Cybernetisk Selskab krever:}
        	\begin{itemize}
        		\item{At du er student og har betalt semesteravgift for inneværende semester.}
        		\item{At du har betalt semesterkontigent til Cybernetisk Selskab.}
        	\end{itemize}
        	\item Medlemmer forplikter seg til å handle i tråd med foreningens formål jf. §1 og etiske retningslinjer for forenings- og styremedlemmer ved instituttet.
        \end{enumerate}
        
        Et vedtektsendringsforslag med vesentlig debatt. Endringsforslaget løftet oppfatningen blant studentmassen på IFI om at driften av Cybernetisk Selskab i stor grad er påvirket av at tidligere aktive og styremedlemmer, og at disse har en for stor rolle i driften til foreningen. Noe som hindrer deltagelsen og tilhørigheten blant aktive studenter. Flere medlemmer er bekymret over de økonomiske problemer som foreningen ville oppleve hvis de måtte betale tilbake prisen for livsvarige medlemskap til alle som hadde kjøpt dem, det ble snakket om summer opp til 50-70 tusen. I denne anledningen ble foreningens betydelige gjeld dratt opp. Dette ble imotsatt av at instituttets studentdrevne Fordelingsutvalg hadde muligheten til å støtte opp om foreningens økonomi. Det ble også nevnt at andelen eldre aktive som møter opp på generalforsamling pleier å være lav, slik det var på nåværende generalforsamling. 
        
        Videre handlet debatten om hvorvidt infrastrukturen til foreningen kunne blitt opprettholdt hvis vi til større grad ekskluderte eldre interne, og påpeker at mye erfaring hadde gått tapt på den måten. Videre ble det løftet at mange interne og medlemmer som kan bidra til studentmiljøet på IFI befinner seg i en fase i livet der det er unaturlig å være semester-registert student, som da ikke nødvendigvis ikke faller under den løse definisjonen av en såkalt pamp.
        
        I det korte ble det også drøftet hvorvidt det hadde vært mulig å rekrutere og bygge opp et nytt frivilliggrunnlag hvis man gjorde slutt med det gamle, det ble løftet fram at andre foreninger klarer dette til stor grad. I denne anledning ble det drøftet at foreningens IT-systemer i stor grad er avhengig av tech-support av den eldre garden, og at tiden var inne for at disse systemene fikk en vårrengjøring i form av et generasjonsskifte. 
        
        Diskutert ble også hvordan det skulle være mulig å være et skjenkested som både var tilgjengelig for offentligheten samtidig som utelukkende studenter kan være medlemmer, hvor forslaget innebar å opprette en database der studenter som hadde betalt inngangspris kunne registrere seg. Blant enkelte var det festet stor tillit til IFI-studenters vilje og evne til å lage organisasjon og systemer fra, om ikke med, scratch.
        
        Få stilte seg imidlertid imot at det kunne være bra for foreningen med en bred debatt om hvordan aktive IFI-studenter kunne føle seg mer involvert i foreningens drift. Andre motstilte seg imidlertid at et vedtektsendringsforslag var riktig måte å få fram en kulturendring på.
        
        Dette viser, med større eller mindre avsporinger, hvordan debatten foregikk. Til slutt ble voteringen tatt to ganger, der sluttresultatet ble at blant 72 gjenværende var 3 ikke tilstede, 44 stemte for, og 25 var enten avholdende eller stemte imot.
        
        \textbf{Forslaget falt}
        

\section{Eventuelt}
    \subsection{Bufferperiode}
        Det foreslås at ved en eventuell vedtektsendring av paragraf 2 at generalforsamlingen vedtar en bufferperiode for når endringen trer i kraft i praksis. Dette for å få utbedret nye rutiner (hvorav registrering i internsystemet, kassesystemet, regnskapet m.m.), gjøre de endringene som trengs i internsystemet og for å få en smidigere overgang ved fjerning av flere medlemskapstyper. Det inngår da at man i situasjoner hvor man må skille mellom medlemmer og ikke-medlemmer, slik som generalforsamling, sjekker studentstatus.
        
        Det foreslås å sette denne bufferperioden til å vare fram til neste semester, både for å spare ekstra arbeid med den umulige kartleggingen av studentstatusen til de 847 medlemmene vi har per dags dato, men også for å gi de som må gjøre alt arbeidet tid til både studiet og eksamener.
        
        \textbf{Forslaget ble ikke realitetsbehandlet, da endringsforslaget under punkt 8.5 falt.}
    
    \subsection{Tilbakemelding på forslagsskjemaet}
        Det ønskes tilbakemelding på skjemaet for vedtektsendringer som ble prøvd ut:
        \href{skjematilbakemelding.cyb.no}{skjematilbakemelding.cyb.no}
        
        De som har brukt skjemaet, samt det tilhørende regnearket, ble oppfordret til å sende inn tilbakemelding på hvordan det var å bruke.
    
    \subsection{Øvrig eventuelt}
    Det takkes for et godt møte. Forsamlingen hyllet debatten.


\part*{Vedlegg:}\label{lastpage}

    \centering\huge Vedtektene kan finnes på \href{vedtekter.cyb.no}{vedtekter.cyb.no}
    
% \addcontentsline{toc}{part}{Vedlegg}

% \newpage
% \phantomsection{}
% \addcontentsline{toc}{section}{Vedtekter for Cybernetisk Selskab} % chktex
% \includepdf[pages=-]{cyb_vedtekter.pdf}

\end{document}
